\subsection{What the source observations cannot guarantee}
\label{sec:v15-identifiability}
A disparity between source-optimal covariance fidelity and source-optimal
held-source AUC does not by itself identify the sign of the
\emph{future} participant's discrimination difference. We formalise
this issue rather than proposing a target-performance predictor from
the observed source-regret ratio.

\paragraph{Proposition 1 (failure of universal source-only dominance).}
Let $w_D$ and $w_C$ be two nonzero discriminant directions learned from
\emph{the same fixed source cohort}. Fix any common target within-class
covariance $\Sigma\succ0$. Define their $\Sigma$-normalised difference
\[
 a=\frac{w_D}{\sqrt{w_D^\top\Sigma w_D}}
   -\frac{w_C}{\sqrt{w_C^\top\Sigma w_C}}.
\]
If $a\ne0$, there exist two binary equal-covariance Gaussian target
populations with \emph{identical covariance and unchanged source data}
for which the relative target AUC of the two discriminants has opposite
signs. Consequently, an observed source-only regularisation-objective
gap cannot uniformly certify which direction improves target AUC
without assumptions restricting the unseen target distribution.

\emph{Proof.}
For a Gaussian target with class-mean contrast $d_g$ and within-class
covariance $\Sigma$, the AUC of direction $w$ is
$\Phi\!\left(w^\top d_g/\sqrt{2w^\top\Sigma w}\right)$.
Because $\Phi$ is increasing, the difference in standardised
separation between $w_D$ and $w_C$ has sign
$\operatorname{sign}(a^\top d_g)$. The two contrasts
$d_g^{(+)}=a$ and $d_g^{(-)}=-a$ give $a^\top d_g^{(+)}
=\|a\|^2>0$ and $a^\top d_g^{(-)}=-\|a\|^2<0$.
Both are valid Gaussian class-mean differences under the same
positive definite covariance. If $a=0$, the directions coincide
after positive $\Sigma$-normalisation and the corresponding
Gaussian AUCs are identical for every $d_g$.\hfill$\square$

\paragraph{Proposition 2 (exact conditional contrast-shift certificate).}
Suppose, \emph{additionally}, that the evaluation population has exactly
the declared shared covariance $\Sigma$ and its unknown class contrast
belongs to the ellipsoid
\[
 \mathcal D_\rho=
 \{d:\,(d-d_s)^\top\Sigma^{-1}(d-d_s)\le\rho^2\},
\]
where $d_s$ is a source-only anchor. The worst possible signed
standardised-margin advantage of $w_D$ over $w_C$ is
\[
 \boxed{
 \inf_{d\in\mathcal D_\rho}
 a^\top d
 =
 a^\top d_s-\rho\sqrt{a^\top\Sigma a}.
 }
\]
If this bound is strictly positive, $w_D$ has greater AUC than
$w_C$ for \emph{every Gaussian evaluation population satisfying
these stated assumptions}. If $a\ne0$, the lower bound is attained
by $d^*=d_s-\rho\,\Sigma a/\sqrt{a^\top\Sigma a}$.
The critical ellipsoid radius is
$\rho^*=(a^\top d_s)/\sqrt{a^\top\Sigma a}$ when
$a^\top d_s>0$.

\emph{Proof.}
Writing $d=d_s+\Sigma^{1/2}z$ with $\|z\|_2\le\rho$,
Cauchy--Schwarz gives
$\min_z a^\top\Sigma^{1/2}z=-\rho\|\Sigma^{1/2}a\|_2$,
with equality at the supplied witness.\hfill$\square$

This result is an \emph{assumption-conditional sensitivity
identity}, not a bound on an arbitrary EEG target population.
In particular, strict Z0 does not let us check a future target's
class contrast or its membership in $\mathcal D_\rho$.
Estimating the radius solely from the observed source cohort
does not establish future coverage. Allowing an unknown covariance
shift or non-Gaussian class distributions additionally invalidates
this simple AUC equivalence. Our implementation therefore emits a
\texttt{CONDITIONAL\_GAUSSIAN\_MATH\_NOT\_EEG\_SAFETY} status
and is kept as a falsification/sensitivity tool rather than a
promoted online policy. A deterministic verification compares
the formula and constructive worst-case witness over 500 randomly
generated positive definite covariances and admissible contrasts.
No new target-subject decoding results come from that test.

\subsection{Four cohorts do not establish universal positive transfer}
\label{sec:v15-four-cohorts}
For a separate, explicitly \emph{post-outcome} external evidence
check, we retained the mean paired discrimination-versus-covariance
contrast of each frozen cohort as one descriptive unit.
The mean AUC differences are $+0.05607$ (EEGMMIDB,
$N=20$), $+0.01763$ (BNCI2014\_002, $N=14$),
$+0.04329$ (BNCI2015\_001, $N=12$), and
$+0.02222$ (Shin2017A, $N=29$). The unweighted
\emph{dataset-mean} effect is $+0.03480$ AUC, with
leave-one-dataset-out means ranging from $+0.02771$ to
$+0.04053$ AUC. Weighting each trial or each overlapping
outer fold as an independent deployment replication would
be inappropriate.

Four positive dataset means would yield a two-sided
$p=0.125$ \emph{if} four independent datasets were
an exchangeable random sample with symmetric null
effect directions. They were \emph{not} drawn under that
sampling design, so this number is a diagnostic
rather than a valid confirmation of improvement in
the population of EEG databases. It cannot replace,
combine or retroactively strengthen the original
separate Holm-controlled test families. The
source-training-pool and participant-harm limitations
in the preceding sections remain in force.

The two propositions explain why a source-only mismatch
can be real and repeatable while providing no
unconditional safe switch for a new individual.
The empirical cross-cohort pattern is
\emph{consistent with} the mechanism but does
not resolve its distributional identifiability gap.
